% Proof fragment for integration into manuscript III; no standalone PDF.
\subsection{最小関係数が強制する真の高次関係}

以下では $A=A_0$, $\mathfrak m=\mathfrak m_A$ と略記し、
$B=A[[s,u,v]]$ と書く。ここで $s=t_{15}$, $u=t_5-11t_{15}$,
$v=t_{16}$ である。加群の Betti 数は
$\beta_i^R(N)=\dim_\kk\operatorname{Tor}_i^R(\kk,N)$ とする。

\begin{lemma}[二つの基底についての最小関係数]\label{lem:v24-betti-comparison}
本稿の加群 $M$ と $Q$ について
\[
 \beta_1^B(M)=12,\qquad
 \beta_1^B(Q)=
 \begin{cases}12&\beta=0,\\13&\beta\ne0\end{cases}
\]
である。一方、
\[
 \beta_1^A(M)\ge4,\qquad
 \beta_1^A(Q)\ge4,
 \qquad \beta\ne0\Longrightarrow\beta_1^A(Q)\ge5.
\]
\end{lemma}

\begin{proof}
$W_9$ を第一列に取った Schur 表示を $E=[b\ C]$ とし、
$C:B^{12}\to B^3$, $M=\operatorname{coker}C$ とする。
四次まで持ち上がる零定数中心 class の空間は
$E_4=\kk[W_9]$ である。もし $Cc=0$ なら $(0,c)$ は $E$ の全次数の核元であり、
その中心値は $E_4$ に入る。第一座標が零なので中心値は零である。
従って $\ker C\subset\mathfrak m_BB^{12}$ である。
さらに $C(0)=0$ であるから、この表示は第一関係まで最小であり、
$\beta_1^B(M)=12$ を得る。

$\beta\ne0$ と仮定する。$E(a,c)=0$ の中心値は同じく $\kk[W_9]$ に入る。
$a$ が単元なら $b=-Ca^{-1}c$ となり $\beta=[b]=0$ に反する。
従って全ての核元の中心値は零で、$E$ は十三本の最小関係を与える。
$\beta=0$ なら $Q=M$ なので、いずれの場合も表示された Betti 数を得る。

閉 $A$-ファイバー $M_0=M/\mathfrak mM$ は五次元で、基底を
$g_0,g_1,g_2,h_3,h_4$ と取ると、非零の作用は
\[
 sg_0=h_3,\qquad sh_3=h_4,\qquad ug_1=h_4,\qquad vg_2=h_4
\]
だけである。三作用についての Koszul 複体の鎖群次元は
$(5,15,15,5)$、微分階数は $(2,5,4)$ である。従って
\[
 \dim_\kk H_1(s,u,v;M_0)=15-2-5=8.
\]
$\beta\in\mathfrak m^2M$ なので $Q/\mathfrak mQ=M_0$ も成立する。

$N=M,Q$ に対して、$A$ の十三パラメータと $s,u,v$ の Koszul 二重複体を用いる。
そのスペクトル系列は
\[
 E^2_{p,q}=H_p\!\left(s,u,v;
       \operatorname{Tor}_q^A(\kk,N)\right)
 \ \Longrightarrow\ \operatorname{Tor}_{p+q}^B(\kk,N)
\]
である。全次数一の二項について
\[
 \dim_\kk E^2_{1,0}=8,\qquad
 \dim_\kk E^2_{0,1}\le\beta_1^A(N)
\]
だから、フィルトレーションの次元を比較して
\[
 \beta_1^B(N)\le8+\beta_1^A(N)
\]
を得る。これと先の十二または十三という値を合わせればよい。
この議論はスペクトル系列の退化を仮定しない。
\end{proof}

\begin{proposition}[循環 torsion の余次元は少なくとも四]\label{prop:v24-height-four}
循環分解
\[
 M\simeq A^4\oplus A/J_M,\qquad
 Q\simeq A^4\oplus A/J_Q,
 \quad J_Q=J_M+(b_0)
\]
におけるイデアルは
\[
 \mu_A(J_M)\ge4,\qquad\mu_A(J_Q)\ge4,
 \qquad \operatorname{ht}J_M,\operatorname{ht}J_Q\ge4
\]
を満たす。特に、三つの独立一次関係を形式座標変換で消去した後にも、
十変数正則局所環上で非零の高次関係イデアルが必ず残る。
$\beta\ne0$ なら $\mu_A(J_Q)\ge5$ である。
\end{proposition}

\begin{proof}
循環分解より
\[
 \beta_1^A(M)=\beta_1^A(A/J_M)=\mu_A(J_M),\qquad
 \beta_1^A(Q)=\mu_A(J_Q).
\]
したがって最小生成数の主張は補題\ref{lem:v24-betti-comparison}による。

$J=J_M$ または $J_Q$ とする。既知の一次項は独立な三形式
\[
 \ell_1=t_8+9t_9-4t_{10}+5t_{14},\qquad
 \ell_2=t_{11}-t_{14},\qquad \ell_3=t_{13}-t_{14}
\]
である。それぞれを一次項に持つ $g_1,g_2,g_3\in J$ を選ぶ。
これらは $A$ の正則パラメータ系の一部であり、
$A'=A/(g_1,g_2,g_3)$ は十変数の形式冪級数環である。
$\mu_A(J)\ge4$ より $J\ne(g_1,g_2,g_3)$ であるため、
像 $J'=J/(g_1,g_2,g_3)$ は $A'$ の非零の真のイデアルである。
$A'$ は整域なので $\operatorname{ht}_{A'}J'\ge1$、従って
\[
 \operatorname{ht}_A J=3+\operatorname{ht}_{A'}J'\ge4.
\]
全ての一次項は消去済みなので $J'\subset\mathfrak m_{A'}^2$ でもある。
これは高次関係の存在を保証するが、その最低次数や生成元を特定するものではない。
\end{proof}

